\documentclass[11pt]{article}
\usepackage[margin=1in]{geometry}
\usepackage{amsmath,amssymb,mathtools}
\usepackage{fontspec,unicode-math}
\setmainfont{Latin Modern Roman}
\setsansfont{Latin Modern Sans}
\setmathfont{Latin Modern Math}
\usepackage{xurl,hyperref}
\hypersetup{colorlinks=true,urlcolor=blue,linkcolor=blue}
\newcommand{\sref}[2]{\hyperlink{source:#1:#2}{[S#2]}}
\newcommand{\cataloguescope}[1]{\paragraph{Recorded mathematical question.}#1}
\begin{document}
\section{A. Navas: A Conjecture on Delone Sets BL to Lattices (after P. Alestalo, D.A. Trotsenko and J. V\"ais\"al\"a). — Problem}\label{statement}

\subsection{Definitions and mathematical statement}

\paragraph{Shared notation.}$\mathbb N=\{1,2,\ldots\}$, and $\mathbb Z,\mathbb Q,\mathbb R,\mathbb C$ denote the integers, rationals, reals and complex numbers. Any other conventions are specified locally.


A Delone set $\mathcal D\subset\mathbb R^2$ is uniformly discrete (distinct points are separated by some $r>0$) and relatively dense (every point of the plane lies within a fixed finite distance of $\mathcal D$). It is bi-Lipschitz equivalent (BL) to $\mathbb Z^2$ if there is a bijection $b:\mathcal D\to\mathbb Z^2$ and $C\geq1$ such that $C^{-1}|x-y|\leq|b(x)-b(y)|\leq C|x-y|$ for all $x,y\in\mathcal D$. A bi-Lipschitz ambient map satisfies the same two-sided bound for all points of the plane.
\paragraph{Statement recorded in the supplied source.}
If $\mathcal D$ is BL to $\mathbb Z^2$, must there exist a bi-Lipschitz map $L:\mathbb R^2\to\mathbb R^2$ such that $L(\mathcal D)=\mathbb Z^2$? \sref{statement}{2}
\subsection{Short English statement}
Can a bi-Lipschitz correspondence between a planar Delone set and the lattice be realized by a bi-Lipschitz map of the whole plane?
\subsection{Sources}
\begin{description}
\item[\hypertarget{source:statement:1}{S1}] ULAM.AI and the UnsolvedMath Contributors, \emph{UnsolvedMath: A Curated Collection of Open Mathematics Problems}, catalogue entry. CC BY 4.0. \url{https://huggingface.co/datasets/ulamai/UnsolvedMath}.
\item[\hypertarget{source:statement:2}{S2}] Faustin Adiceam (compiler), \emph{Open Problems and Conjectures related to the Theory of Mathematical Quasicrystals}, arXiv:1604.06280v2 (2016), Section 2.6, Problem 2.6.1, pp. 8--9. \url{https://arxiv.org/pdf/1604.06280}.
\end{description}
\label{end:statement}
\end{document}
